\subsubsection{Subretícula invariante y unidad cuadrática del operador agregado}
\label{mg02:subreticula}

La agregación por clases de residuos conserva una estructura integral que
precisa la descomposición espectral anterior. El operador de partida es
\(M_\Pi\), obtenido al promediar la hoja multiplicativa de APP sobre
los bloques \(C_a\times C_b\). Su dominio integral es \(\mathbb Z^3\),
con las coordenadas ordenadas por \(C_1,C_2,C_0\).

\begin{proposition}[Descomposición sobre una subretícula de índice dos]
Sean
\[
 v_-=(1,-1,0)^{\mathsf T},\qquad
 v_+=(1,1,0)^{\mathsf T},\qquad v_0=(0,0,1)^{\mathsf T},
\]
y sea \(L=\mathbb Zv_-\oplus\mathbb Zv_+\oplus\mathbb Zv_0\).
Entonces
\begin{equation}
 L=\{(x,y,z)\in\mathbb Z^3:x\equiv y\pmod2\},
 \qquad [\mathbb Z^3:L]=2.
 \label{mg02:indice}
\end{equation}
La subretícula \(L\) es invariante por \(M_\Pi\). En la base
ordenada \((v_-,v_+,v_0)\), la restricción tiene la matriz
\begin{equation}
 [M_\Pi|_L]=(-1)\oplus3H,\qquad
 H=\begin{pmatrix}3&2\\4&3\end{pmatrix},\qquad \det H=1.
 \label{mg02:restriccion}
\end{equation}
\end{proposition}

\begin{proof}
La combinación \(av_-+bv_++cv_0=(a+b,-a+b,c)\) tiene las dos
primeras coordenadas de igual paridad. Recíprocamente, para un vector
con esa paridad, las coordenadas en la base indicada son
\(a=(x-y)/2\), \(b=(x+y)/2\) y \(c=z\), todas enteras. La matriz
de cambio de base y su inversa racional son
\[
 P=\begin{pmatrix}1&1&0\\-1&1&0\\0&0&1\end{pmatrix},\qquad
 P^{-1}=\begin{pmatrix}1/2&-1/2&0\\1/2&1/2&0\\0&0&1\end{pmatrix}.
\]
En este orden de columnas, \(\det P=2\), que da el índice de
\eqref{mg02:indice}. La acción sobre los tres generadores es
\[
 M_\Pi v_-=-v_-,\qquad
 M_\Pi v_+=9v_++12v_0,\qquad
 M_\Pi v_0=6v_++9v_0.
\]
Estas igualdades prueban a la vez la invariancia y la identidad
\[
 P^{-1}M_\Pi P=
 \begin{pmatrix}-1&0&0\\0&9&6\\0&12&9\end{pmatrix}
 =(-1)\oplus3H.
\]
Por último, \(\det H=9-8=1\).
\end{proof}

\begin{corollary}[Iteración de Pell en el submódulo de rango dos]
Para \(n\geq0\) existen enteros \(a_n,b_n\) tales que
\begin{equation}
 H^n=\begin{pmatrix}a_n&b_n\\2b_n&a_n\end{pmatrix},\qquad
 a_n+b_n\sqrt2=(3+2\sqrt2)^n,
 \qquad a_n^2-2b_n^2=1.
 \label{mg02:pell}
\end{equation}
En particular, \((a_0,b_0)=(1,0)\), \((a_1,b_1)=(3,2)\) y
\((a_2,b_2)=(17,12)\). Los autovalores de \(H\) son las unidades
cuadráticas conjugadas \(3\pm2\sqrt2\).
\end{corollary}

\begin{proof}
La multiplicación de matrices de la forma indicada corresponde a
\[
 (a,b)(c,d)=(ac+2bd,ad+bc),
\]
que es la ley de multiplicación de \(a+b\sqrt2\). Así, las potencias
de \(H\) satisfacen
\(a_{n+1}=3a_n+4b_n\), \(b_{n+1}=2a_n+3b_n\), con los datos
iniciales anteriores. La norma multiplicativa de \(3+2\sqrt2\) vale
uno; por tanto, \(a_n^2-2b_n^2=1\) en todo grado. El polinomio
característico de \(H\), \(\lambda^2-6\lambda+1\), da las dos
unidades conjugadas. La identidad
\(H^{\mathsf T}\operatorname{diag}(2,-1)H=
\operatorname{diag}(2,-1)\) expresa la misma conservación como
una forma bilineal integral.
\end{proof}

La separación integral tiene, por tanto, residencia en \(L\); la
misma matriz \(P\) da una descomposición completa sobre
\(\mathbb R^3\). El factor \(3\) de la restricción conserva la
escala del promedio APP: su iteración de rango dos es \(3^nH^n\).
El bloque central \(B_c=\left(\begin{smallmatrix}7&2\\2&7\end{smallmatrix}\right)\)
reside, en cambio, en dos posiciones adyacentes de la tabla y satisface
\(B_c^{\mathsf T}\operatorname{diag}(1,-1)B_c=
45\operatorname{diag}(1,-1)\). La compresión en tres clases y la
restricción central en dos posiciones conservan así sus dominios,
bases y normalizaciones respectivos. La estructura de Pell procede de
la primera operación mediante \eqref{mg02:restriccion}.
